% =====================================================================
%  Conditional Things in ML for Mathematicians
%  A handbook on the rigorous meaning of conditional notation in ML.
%  Compile:  pdflatex main.tex  (run twice for TOC/refs)
% =====================================================================
\documentclass[11pt,oneside]{report}

% ---- page geometry & fonts ----
\usepackage[a4paper,margin=1in]{geometry}
\usepackage[T1]{fontenc}
\usepackage{lmodern}
\usepackage{microtype}

% ---- math ----
\usepackage{amsmath,amssymb,amsthm}
\usepackage{mathtools}      % \DeclarePairedDelimiter, \coloneqq, etc.
\usepackage{mathrsfs}       % \mathscr
\usepackage{bm}

% ---- layout / misc ----
\usepackage{enumitem}
\usepackage{booktabs}
\usepackage{longtable}
\usepackage{array}
\usepackage{xcolor}
\usepackage[bottom]{footmisc}
\usepackage{emptypage}

% ---- hyperref must be (nearly) last ----
\usepackage[colorlinks=true,linkcolor=black!75!blue,citecolor=black,
            urlcolor=blue!60!black,linktoc=all]{hyperref}

% =====================================================================
%  Theorem environments
% =====================================================================
\theoremstyle{plain}
\newtheorem{theorem}{Theorem}[chapter]
\newtheorem{proposition}[theorem]{Proposition}
\newtheorem{lemma}[theorem]{Lemma}
\newtheorem{corollary}[theorem]{Corollary}

\theoremstyle{definition}
\newtheorem{definition}[theorem]{Definition}
\newtheorem{notation}[theorem]{Notation}
\newtheorem{convention}[theorem]{Convention}
\newtheorem{example}[theorem]{Example}

\theoremstyle{remark}
\newtheorem{remark}[theorem]{Remark}
\newtheorem{warning}[theorem]{Warning}
\newtheorem*{warning*}{\textbf{Pitfall}}

% A lightweight environment for the symbol-dictionary entries (Part VI).
\newenvironment{mlentry}[1]%
  {\par\medskip\noindent\textbf{Symbol.}\quad #1\par\nobreak\smallskip}%
  {\par\medskip\hrule height0pt\medskip}

% =====================================================================
%  Macros
% =====================================================================
\newcommand{\R}{\mathbb{R}}
\newcommand{\N}{\mathbb{N}}
\newcommand{\Z}{\mathbb{Z}}
\newcommand{\Prob}{\mathbb{P}}
\newcommand{\E}{\mathbb{E}}
\newcommand{\one}{\mathbf{1}}                 % indicator
\newcommand{\dd}{\,\mathrm{d}}
\newcommand{\Fcal}{\mathcal{F}}
\newcommand{\Gcal}{\mathcal{G}}
\newcommand{\Hcal}{\mathcal{H}}
\newcommand{\Bcal}{\mathcal{B}}
\newcommand{\Xcal}{\mathcal{X}}
\newcommand{\Ycal}{\mathcal{Y}}
\newcommand{\Zcal}{\mathcal{Z}}
\newcommand{\Dcal}{\mathcal{D}}
\newcommand{\eqas}{\overset{\text{a.s.}}{=}}
\newcommand{\defeq}{\vcentcolon=}
\newcommand{\given}{\,\vert\,}
\newcommand{\biggiven}{\,\middle\vert\,}

% independence / conditional independence
\newcommand{\indep}{\mathrel{\perp\mkern-10mu\perp}}

% operators
\DeclareMathOperator{\Var}{Var}
\DeclareMathOperator{\Cov}{Cov}
\DeclareMathOperator{\Cor}{Cor}
\DeclareMathOperator{\KL}{KL}
\DeclareMathOperator{\supp}{supp}
\DeclareMathOperator{\dom}{dom}
\DeclareMathOperator{\sgn}{sgn}
\DeclareMathOperator*{\argmin}{arg\,min}
\DeclareMathOperator*{\argmax}{arg\,max}
\DeclareMathOperator{\Ent}{H}        % entropy
\DeclareMathOperator{\MI}{I}         % mutual information
\DeclareMathOperator{\diffent}{h}    % differential entropy

% paired delimiters
\DeclarePairedDelimiter{\abs}{\lvert}{\rvert}
\DeclarePairedDelimiter{\norm}{\lVert}{\rVert}
\DeclarePairedDelimiter{\inner}{\langle}{\rangle}

% conditional expectation with a nice scalable bar:  \cE{X}{\Gcal}
\newcommand{\cE}[2]{\E\!\left[\,#1 \,\middle|\, #2\,\right]}
\newcommand{\cProb}[2]{\Prob\!\left(\,#1 \,\middle|\, #2\,\right)}
\newcommand{\cVar}[2]{\Var\!\left(\,#1 \,\middle|\, #2\,\right)}

\allowdisplaybreaks
\numberwithin{equation}{chapter}

% =====================================================================
\begin{document}

% ---------------------------------------------------------------------
%  Title page
% ---------------------------------------------------------------------
\begin{titlepage}
\centering
\vspace*{3.5cm}
{\Huge\bfseries Conditional Things in ML\\[0.35em] for Mathematicians\par}
\vspace{1.2cm}
{\Large A Rigorous Handbook of the Conditional Notation\\
Used in Machine-Learning Theses\par}
\vfill
{\large A measure-theoretic reading of $\;\Prob(A\given B),\quad
\cE{X}{\Gcal},\quad p(x\given y),\quad X\indep Y\given Z,\quad
H(X\given Y),\quad \nabla_x\log p_t(x\given y).$\par}
\vfill
{\large From conditional probability to classifier-free guidance,\\
with every bar ``$\,\vert\,$'' given a definition.\par}
\vspace{2cm}
{\large\today\par}
\end{titlepage}

\pagenumbering{roman}
\tableofcontents

% ---------------------------------------------------------------------
%  Preface
% ---------------------------------------------------------------------
\chapter*{Preface}
\addcontentsline{toc}{chapter}{Preface}
\markboth{Preface}{Preface}

Machine-learning papers are written in a dialect of probability theory in
which the vertical bar ``$\,\vert\,$'' is the busiest symbol on the page.
A single thesis may, within three paragraphs, write the conditional
density $p(x\given y)$, the conditional expectation $\E[X\given Y]$, the
amortized variational posterior $q_\phi(z\given x)$, the policy
$\pi(a\given s)$, and the conditional score $\nabla_x\log p_t(x\given y)$
--- and treat all of them as if their meaning were self-evident. For a
reader trained in measure theory this is disorienting: the same glyph
denotes a number, a function, a random variable, a Markov kernel, and a
Radon--Nikodym derivative, depending on context that is rarely stated.

This handbook is a translation manual. Its purpose is to take every
common piece of \emph{conditional} notation that a mathematician will
meet in an ML thesis and attach to it a single, unambiguous,
measure-theoretic definition, together with the theorem that guarantees
the object exists, the sense in which it is unique, and the pitfalls that
the informal notation hides.

\paragraph{Who this is for.}
The intended reader knows graduate real analysis and measure-theoretic
probability at the level of, say, the first chapters of Williams'
\emph{Probability with Martingales} or Durrett's \emph{Probability:
Theory and Examples}: $\sigma$-algebras, the Lebesgue integral, product
measures, Radon--Nikodym, and the basic convergence theorems are assumed.
No prior machine learning is assumed. Where an ML construction (a VAE, a
diffusion model, an MDP) is needed, it is introduced from its probabilistic
skeleton only, with enough detail to fix the notation and not one symbol more.

\paragraph{What this is not.}
It is not a machine-learning textbook, an optimization reference, or a
survey of architectures. Approximation, estimation, and computation ---
the parts of ML that are genuinely about \emph{learning} --- appear only
insofar as they are needed to explain what a symbol denotes. We are
concerned with what the notation \emph{means}, not with how well a network
fits it.

\paragraph{Standing conventions.}
Unless stated otherwise, $(\Omega,\Fcal,\Prob)$ is a fixed probability
space; random variables are measurable maps out of it; ``a.s.'' means
$\Prob$-almost surely; and all named random variables are assumed to take
values in standard Borel (Polish, with their Borel $\sigma$-algebra)
spaces, which is exactly the hypothesis that makes regular conditional
distributions exist (Chapter~\ref{ch:cond-dist}). Densities are taken with
respect to a reference measure --- Lebesgue on $\R^d$ or counting measure
on a discrete set --- that we name when it matters and suppress when it
does not. A short index of recurring symbols is collected in
Appendix~\ref{app:cheatsheet}.

\paragraph{How to read it.}
Parts~\ref{part:foundations}--\ref{part:info} build the rigorous objects
in order of dependence: conditional expectation, then conditional
distributions and densities, then conditional independence, then the
information-theoretic conditionals. Part~\ref{part:dictionary} is a
dictionary: each entry takes a symbol \emph{as it is actually printed in
papers} and points back to the definition that licenses it. A reader who
only wants to decode a particular formula can start there and follow the
cross-references inward.

\vspace{1em}
\noindent\emph{A note on rigor and honesty.} Several identities that ML
papers write with an equals sign are, strictly, only true almost surely,
or only after a regular version has been fixed, or only under an absolute-continuity
hypothesis that is left implicit. Rather than smooth these over, the
handbook flags them, because the gap between ``$=$'' and ``$\eqas$'' is
exactly where a mathematician's intuition is most likely to be misled.

\cleardoublepage
\pagenumbering{arabic}

% =====================================================================
\part{Foundations and Conditional Expectation}\label{part:foundations}
% =====================================================================

\chapter{Why elementary conditioning is not enough}\label{ch:why}

\section{The schoolbook definition and where it breaks}

Every reader knows the elementary definition.

\begin{definition}[Elementary conditional probability]\label{def:elem-cond}
Let $(\Omega,\Fcal,\Prob)$ be a probability space and $B\in\Fcal$ an event
with $\Prob(B)>0$. For $A\in\Fcal$,
\begin{equation}\label{eq:elem-cond}
  \cProb{A}{B}\;\defeq\;\frac{\Prob(A\cap B)}{\Prob(B)}.
\end{equation}
\end{definition}

For fixed $B$ with $\Prob(B)>0$ the set function $A\mapsto\cProb{A}{B}$ is
a bona fide probability measure on $(\Omega,\Fcal)$, and
$\cE{X}{B}\defeq\Prob(B)^{-1}\int_B X\dd\Prob$ is its expectation operator.
Nothing here is subtle. The trouble is that essentially every conditional
quantity in machine learning conditions on an event of probability zero.

\begin{example}[The event we actually want has probability zero]
\label{ex:zero}
Let $(X,Y)$ be jointly continuous on $\R^2$, e.g.\ $Y\sim\mathcal N(0,1)$
and $X\given Y$ Gaussian. The object an ML paper denotes $p(x\given y)$,
$\cE{X}{Y=y}$, or $\Prob(X\in A\given Y=y)$ conditions on the event
$\{Y=y\}$, which has $\Prob(Y=y)=0$. Formula~\eqref{eq:elem-cond} returns
$\tfrac{0}{0}$. The symbol is in universal use and the schoolbook
definition cannot assign it a value.
\end{example}

The naive repair --- ``condition on $\{y-\varepsilon<Y<y+\varepsilon\}$ and
let $\varepsilon\to0$'' --- can be made to work, but its answer depends on
\emph{how} the limiting family of events is chosen. That dependence is not
a technicality; it is a genuine ambiguity with a name.

\begin{warning*}[The Borel--Kolmogorov paradox, previewed]
Conditioning on a measure-zero event $\{Y=y\}$ has no answer until one
specifies the \emph{family} of shrinking events used to approach it.
Different parameterizations of the same null event yield different
``conditional distributions.'' We treat this carefully in
Section~\ref{sec:bk-paradox}; for now it is the warning that motivates
everything below. The resolution is to stop conditioning on a single
event and instead condition on a $\sigma$-algebra, obtaining an object
defined only up to almost-sure equivalence.
\end{warning*}

\section{The shift in viewpoint}

Measure-theoretic conditioning makes three moves at once, and a
mathematician new to the ML literature should hold all three in mind.

\begin{enumerate}[label=(\roman*),leftmargin=2.2em]
\item \textbf{Condition on information, not on an event.} The primitive is
a sub-$\sigma$-algebra $\Gcal\subseteq\Fcal$ --- ``everything we are
allowed to know.'' Conditioning on a random variable $Y$ means
conditioning on $\Gcal=\sigma(Y)$.
\item \textbf{The answer is a random variable, not a number.}
$\cE{X}{\Gcal}$ is itself a $\Gcal$-measurable random variable; evaluating
it ``at $Y=y$'' is a secondary, and more delicate, operation
(Section~\ref{sec:cond-on-rv}).
\item \textbf{Identities hold almost surely.} The defining property pins
down $\cE{X}{\Gcal}$ only up to a $\Prob$-null set. Every ``$=$'' between
conditional objects is, by default, an ``$\eqas$''.
\end{enumerate}

The remainder of Part~\ref{part:foundations} develops move (i)--(iii) into
Kolmogorov's definition of conditional expectation, from which conditional
probability, conditional variance, and (in Part~\ref{part:cond-dist}) the
conditional densities and kernels of the ML literature are all derived as
special cases.

% =====================================================================
\part{Conditional Expectation}\label{part:cond-exp}
% =====================================================================

\chapter{Conditional expectation given a \texorpdfstring{$\sigma$}{sigma}-algebra}
\label{ch:cond-exp}

\section{Kolmogorov's definition}

\begin{definition}[Conditional expectation]\label{def:cond-exp}
Let $X\in L^1(\Omega,\Fcal,\Prob)$ and let $\Gcal\subseteq\Fcal$ be a
sub-$\sigma$-algebra. A \emph{version of the conditional expectation of
$X$ given $\Gcal$} is any random variable $Z$ satisfying
\begin{enumerate}[label=(CE\arabic*),leftmargin=3.2em]
\item $Z$ is $\Gcal$-measurable;\label{ce:meas}
\item $Z\in L^1(\Prob)$, and for every $G\in\Gcal$,
      \begin{equation}\label{eq:ce-defining}
        \int_G Z\dd\Prob \;=\; \int_G X\dd\Prob .
      \end{equation}
\end{enumerate}
Any such $Z$ is written $\cE{X}{\Gcal}$.
\end{definition}

Property~\ref{ce:meas} says $Z$ uses only the information in $\Gcal$;
property~\eqref{eq:ce-defining} says $Z$ has the same ``local averages'' as
$X$ over every event we can resolve with $\Gcal$. Note what is \emph{not}
required: no pointwise relation between $Z(\omega)$ and $X(\omega)$ holds,
only an averaged one.

\begin{theorem}[Existence and a.s.\ uniqueness]\label{thm:ce-exists}
For every $X\in L^1(\Prob)$ and every sub-$\sigma$-algebra
$\Gcal\subseteq\Fcal$ a version $\cE{X}{\Gcal}$ exists, and any two
versions agree $\Prob$-almost surely.
\end{theorem}

\begin{proof}[Proof sketch]
\emph{Existence.} Assume first $X\ge 0$. Define a finite measure on
$(\Omega,\Gcal)$ by $\nu(G)\defeq\int_G X\dd\Prob$. Then $\nu\ll\Prob|_\Gcal$,
since $\Prob(G)=0\Rightarrow\nu(G)=0$. By the Radon--Nikodym theorem
(Theorem~\ref{thm:rn}) applied on the measurable space $(\Omega,\Gcal)$
there is a $\Gcal$-measurable density $Z=\dd\nu/\dd(\Prob|_\Gcal)$, which is
exactly~\eqref{eq:ce-defining}. For general $X$ set
$\cE{X}{\Gcal}=\cE{X^+}{\Gcal}-\cE{X^-}{\Gcal}$.

\emph{Uniqueness.} If $Z,Z'$ both satisfy~\eqref{eq:ce-defining}, then
$\int_G(Z-Z')\dd\Prob=0$ for all $G\in\Gcal$. Taking
$G=\{Z-Z'>0\}\in\Gcal$ forces $\Prob(Z>Z')=0$, and symmetrically; hence
$Z\eqas Z'$.
\end{proof}

\begin{convention}[``The'' conditional expectation]
Because versions differ only on a null set, one speaks of \emph{the}
conditional expectation and writes $\cE{X}{\Gcal}$ for an arbitrary
version, with all ensuing equalities understood $\Prob$-a.s. When a
statement is made for a \emph{specific} version (as in the regular
conditional distributions of Chapter~\ref{ch:cond-dist}), we say so
explicitly.
\end{convention}

\begin{remark}[$L^2$ projection picture]\label{rem:l2-proj}
If $X\in L^2(\Prob)$ then $L^2(\Omega,\Gcal,\Prob)$ is a closed subspace of
$L^2(\Omega,\Fcal,\Prob)$, and $\cE{X}{\Gcal}$ is exactly the orthogonal
projection of $X$ onto it. Equivalently it solves
\begin{equation}
  \cE{X}{\Gcal}\;=\;\argmin_{\substack{Z\ \Gcal\text{-meas.}\\ Z\in L^2}}
     \;\E\bigl[(X-Z)^2\bigr],
\end{equation}
the \emph{best mean-square predictor of $X$ from the information
$\Gcal$.} This is the identity that ML's ``the regression function
$\E[Y\given X{=}\,\cdot\,]$ minimizes squared error'' is quoting, and it is
the cleanest one-line answer to ``what is conditional expectation?'' for
an analyst: it is a projection.
\end{remark}

\section{The calculus of conditional expectation}\label{sec:ce-calculus}

The following properties are used constantly --- usually silently --- in
ML derivations (the ELBO, policy-gradient identities, the tower rule
behind bootstrapping in reinforcement learning). All equalities are
$\Prob$-a.s.; all random variables are in $L^1$ unless a moment is named.

\begin{proposition}[Elementary properties]\label{prop:ce-props}
Let $\Gcal,\Hcal\subseteq\Fcal$ be sub-$\sigma$-algebras and $a,b\in\R$.
\begin{enumerate}[label=\textup{(\arabic*)},leftmargin=2.4em]
\item \textup{(Linearity)} $\cE{aX+bY}{\Gcal}=a\cE{X}{\Gcal}+b\cE{Y}{\Gcal}$.
\item \textup{(Monotonicity)} $X\le Y$ a.s.\ $\Rightarrow
      \cE{X}{\Gcal}\le\cE{Y}{\Gcal}$; in particular
      $\bigl|\cE{X}{\Gcal}\bigr|\le\cE{\abs{X}}{\Gcal}$.
\item \textup{(Total expectation)} $\E\bigl[\cE{X}{\Gcal}\bigr]=\E[X]$,
      taking $G=\Omega$ in~\eqref{eq:ce-defining}.
\item \textup{(Tower / iterated expectations)} if $\Hcal\subseteq\Gcal$ then
      \begin{equation}\label{eq:tower}
        \cE{\,\cE{X}{\Gcal}\,}{\Hcal}\;=\;\cE{X}{\Hcal}
        \;=\;\cE{\,\cE{X}{\Hcal}\,}{\Gcal}.
      \end{equation}
\item \textup{(Taking out what is known)} if $Y$ is $\Gcal$-measurable and
      $XY\in L^1$ (e.g.\ $Y$ bounded), then
      \begin{equation}\label{eq:take-out}
        \cE{YX}{\Gcal}\;=\;Y\,\cE{X}{\Gcal}.
      \end{equation}
\item \textup{(Independence drops the conditioning)} if $\sigma(X)$ is
      independent of $\Gcal$, then $\cE{X}{\Gcal}=\E[X]$ (a constant).
\item \textup{(Conditional Jensen)} if $\varphi:\R\to\R$ is convex and
      $\varphi(X)\in L^1$, then
      \begin{equation}\label{eq:cond-jensen}
        \varphi\!\bigl(\cE{X}{\Gcal}\bigr)\;\le\;\cE{\varphi(X)}{\Gcal}.
      \end{equation}
      With $\varphi(t)=\abs{t}^p$ ($p\ge1$) this gives the conditional
      $L^p$ contraction $\bigl\|\cE{X}{\Gcal}\bigr\|_p\le\norm{X}_p$.
\end{enumerate}
\end{proposition}

\begin{proof}[Notes on proof]
(1)--(3) are immediate from Definition~\ref{def:cond-exp}. For the tower
rule, the left equality holds because $\cE{X}{\Hcal}$ is $\Hcal$-measurable
and satisfies~\eqref{eq:ce-defining} over $\Hcal\subseteq\Gcal$; the right
equality is~\eqref{eq:ce-defining} read for $\Hcal$. Property~(5) is proved
first for $Y=\one_{G_0}$, $G_0\in\Gcal$ (both sides integrate to
$\int_{G\cap G_0}X$ over $G\in\Gcal$), then extended by linearity and
monotone/dominated convergence. Property~(7) follows from the unconditional
Jensen inequality applied to the probability measures $\cProb{\cdot}{\Gcal}$
of Chapter~\ref{ch:cond-dist}, or directly via suprema of affine minorants
of $\varphi$.
\end{proof}

\begin{remark}[Why ML cares about the tower rule]
The identity~\eqref{eq:tower} is the backbone of \emph{bootstrapping}.
In reinforcement learning the value function satisfies
$V(s)=\E\!\left[r+\gamma V(S')\given S=s\right]$ precisely because a
total-reward conditional expectation can be peeled one step at a time by
the tower rule; see Section~\ref{sec:rl}. In variational inference the same
rule lets one push an expectation through a latent variable to derive the
ELBO (Section~\ref{sec:elbo}).
\end{remark}

\section{Conditional probability and conditional variance}

\begin{definition}[Conditional probability given $\Gcal$]\label{def:cond-prob-sigma}
For $A\in\Fcal$ define the random variable
\begin{equation}
  \cProb{A}{\Gcal}\;\defeq\;\cE{\one_A}{\Gcal}.
\end{equation}
\end{definition}

By Proposition~\ref{prop:ce-props}, for each \emph{fixed} $A$ we have
$0\le\cProb{A}{\Gcal}\le1$ a.s., $\cProb{\Omega}{\Gcal}=1$ a.s., and
countable additivity holds a.s.\ for any fixed countable disjoint family.
The crucial subtlety --- that these ``a.s.'' qualifiers cannot in general
be made simultaneous over \emph{all} $A$ --- is the entire content of
Chapter~\ref{ch:cond-dist}.

\begin{definition}[Conditional variance]\label{def:cond-var}
For $X\in L^2(\Prob)$,
\begin{equation}
  \cVar{X}{\Gcal}\;\defeq\;\cE{\bigl(X-\cE{X}{\Gcal}\bigr)^2}{\Gcal}
   \;=\;\cE{X^2}{\Gcal}-\bigl(\cE{X}{\Gcal}\bigr)^2 .
\end{equation}
\end{definition}

\begin{proposition}[Law of total variance]\label{prop:total-var}
For $X\in L^2(\Prob)$,
\begin{equation}\label{eq:total-var}
  \Var(X)\;=\;\E\bigl[\cVar{X}{\Gcal}\bigr]
            \;+\;\Var\bigl(\cE{X}{\Gcal}\bigr).
\end{equation}
\end{proposition}
\begin{proof}
Expand both terms using the tower rule and
$\E[\cE{X}{\Gcal}]=\E[X]$; the cross term vanishes because
$\cE{X}{\Gcal}$ is the $L^2$-projection of $X$ (Remark~\ref{rem:l2-proj}),
so $X-\cE{X}{\Gcal}\perp L^2(\Gcal)\ni\cE{X}{\Gcal}-\E[X]$.
\end{proof}

\begin{remark}
Equation~\eqref{eq:total-var} is the ``bias/variance'' bookkeeping behind,
e.g., the decomposition of predictive uncertainty into \emph{aleatoric}
($\E[\cVar{X}{\Gcal}]$, irreducible noise) and \emph{epistemic}
($\Var(\cE{X}{\Gcal})$, reducible by more information) parts that recurs
throughout the Bayesian deep-learning literature.
\end{remark}

\chapter{Conditioning on a random variable}\label{ch:cond-on-rv}

\section{From \texorpdfstring{$\sigma(Y)$}{sigma(Y)} to a function of \texorpdfstring{$y$}{y}}
\label{sec:cond-on-rv}

ML papers almost never write $\cE{X}{\Gcal}$; they write $\cE{X}{Y}$ and,
more often, $\cE{X}{Y=y}$. Making sense of the second requires one
structural theorem.

\begin{definition}[Conditioning on a random variable]
For a random variable $Y:\Omega\to\Ycal$ set $\sigma(Y)\defeq
Y^{-1}(\Bcal_\Ycal)$ and define
\(
  \cE{X}{Y}\defeq\cE{X}{\sigma(Y)},\qquad
  \cProb{A}{Y}\defeq\cProb{A}{\sigma(Y)}.
\)
\end{definition}

\begin{lemma}[Doob--Dynkin]\label{lem:doob-dynkin}
Let $Y:\Omega\to\Ycal$ and $Z:\Omega\to\R$ be measurable, with $\Ycal$
standard Borel. Then $Z$ is $\sigma(Y)$-measurable if and only if there
exists a Borel function $g:\Ycal\to\R$ with $Z=g(Y)$.
\end{lemma}

Applied to $Z=\cE{X}{Y}$, which is $\sigma(Y)$-measurable by construction,
the lemma produces a Borel $g$ with $\cE{X}{Y}=g(Y)$ a.s. \emph{This $g$ is
what licenses the notation $\cE{X}{Y=y}$:}

\begin{definition}[The evaluated conditional expectation]\label{def:eval-ce}
With $g$ as above, define
\begin{equation}
  \cE{X}{Y=y}\;\defeq\;g(y),\qquad y\in\Ycal.
\end{equation}
\end{definition}

\begin{warning}[What ``$\cE{X}{Y=y}$'' really is]\label{warn:version}
The function $g$ of Lemma~\ref{lem:doob-dynkin} is determined only up to a
$P_Y$-null set, where $P_Y=\Prob\circ Y^{-1}$ is the law of $Y$. Hence
$\cE{X}{Y=y}$ is well defined only for $P_Y$-\emph{almost every} $y$, and
its value at any single, pre-specified $y_0$ is meaningless in isolation.
Two consequences a mathematician should hold onto:
\begin{itemize}[leftmargin=1.6em]
\item Statements like ``$\cE{X}{Y=y}=h(y)$ \emph{for all} $y$'' are abuses:
they assert equality of two versions, true only $P_Y$-a.e.
\item The random variable $\cE{X}{Y}=g(Y)$ \emph{is} well defined a.s.,
because $Y$ never visits the $P_Y$-null exceptional set. The pathology lives
entirely in the act of plugging in a deterministic $y$.
\end{itemize}
\end{warning}

\begin{theorem}[Change of variables / the law of $\cE{X}{Y}$]\label{thm:cev}
With $g(y)=\cE{X}{Y=y}$, for every bounded Borel $\psi:\Ycal\to\R$,
\begin{equation}\label{eq:cev}
  \E\bigl[X\,\psi(Y)\bigr]\;=\;\E\bigl[g(Y)\,\psi(Y)\bigr]
  \;=\;\int_\Ycal g(y)\,\psi(y)\,P_Y(\dd y).
\end{equation}
\end{theorem}
\begin{proof}
The first equality is~\eqref{eq:take-out} with the $\sigma(Y)$-measurable
factor $\psi(Y)$, then total expectation; the second is the image-measure
(``law of the unconscious statistician'') formula.
\end{proof}

Equation~\eqref{eq:cev} is the rigorous skeleton of the ubiquitous ML move
\(
  \E_{X,Y}[\,\cdot\,]=\E_{Y}\bigl[\E_{X\given Y}[\,\cdot\,]\bigr],
\)
written in papers as ``condition on $Y$, then average.''

\section{The regression function}

\begin{definition}[Regression function]
For $X\in L^1$ the map $m(y)\defeq\cE{X}{Y=y}$ is the \emph{regression
function} of $X$ on $Y$.
\end{definition}

By Remark~\ref{rem:l2-proj}, when $X\in L^2$ the random variable $m(Y)$ is
the $L^2(\sigma(Y))$ projection of $X$, hence
\begin{equation}
  m\;=\;\argmin_{f\ \text{Borel},\,f(Y)\in L^2}\ \E\bigl[(X-f(Y))^2\bigr].
\end{equation}
This is precisely the population object that a regression network with the
squared-error loss is trying to approximate; the Bayes-optimal predictor of
$X$ from $Y$ under squared loss \emph{is} the conditional expectation.
Different losses target different conditional functionals: absolute loss
targets a conditional median, and the ``pinball'' loss at level $\tau$
targets the conditional $\tau$-quantile --- a point worth keeping in mind
when a paper says a network ``predicts $Y$ given $X$'' without naming the
loss.

% =====================================================================
\part{Conditional Distributions, Kernels, and Densities}\label{part:cond-dist}
% =====================================================================

\chapter{Markov kernels and regular conditional distributions}
\label{ch:cond-dist}

Conditional expectation gave us, for each \emph{fixed} event $A$, a random
variable $\cProb{A}{\Gcal}$. What an ML paper means by ``the conditional
distribution $p(\cdot\given y)$'' is the far stronger object in which $A$
and $\omega$ vary together and the result is, for almost every outcome, an
honest probability measure. The bridge is the notion of a Markov kernel.

\section{Markov kernels}

\begin{definition}[Markov kernel]\label{def:kernel}
Let $(\Ycal,\Bcal_\Ycal)$ and $(\Xcal,\Bcal_\Xcal)$ be measurable spaces.
A \emph{Markov kernel} (or \emph{transition kernel}, \emph{stochastic
kernel}, \emph{conditional distribution}) from $\Ycal$ to $\Xcal$ is a map
$\kappa:\Ycal\times\Bcal_\Xcal\to[0,1]$ such that
\begin{enumerate}[label=\textup{(K\arabic*)},leftmargin=3.2em]
\item for each fixed $y\in\Ycal$, $A\mapsto\kappa(y,A)$ is a probability
      measure on $(\Xcal,\Bcal_\Xcal)$;\label{k:meas}
\item for each fixed $A\in\Bcal_\Xcal$, $y\mapsto\kappa(y,A)$ is
      $\Bcal_\Ycal$-measurable.\label{k:meas2}
\end{enumerate}
We write $\kappa(y,\dd x)$ or $\kappa(\dd x\given y)$, and overload
notation by writing $\kappa(y,\cdot)=:p(\cdot\given y)$ once a reference
measure is fixed.
\end{definition}

A Markov kernel is exactly ``a measurably-indexed family of probability
measures.'' It is the rigorous content of every conditional distribution
symbol in the ML literature: $p_\theta(x\given y)$, $q_\phi(z\given x)$,
$\pi(\cdot\given s)$, the forward and reverse transitions of a diffusion,
and the rows of a Markov-chain transition matrix are all Markov kernels.

\begin{remark}[Kernels compose; this is why ``stacking conditionals''
works]\label{rem:kernel-compose}
Given a measure $\mu$ on $\Ycal$ and a kernel $\kappa$ from $\Ycal$ to
$\Xcal$, one builds a measure on the product and a marginal:
\begin{equation}
  (\mu\otimes\kappa)(B\times A)=\int_B \kappa(y,A)\,\mu(\dd y),\qquad
  (\mu\kappa)(A)=\int_\Ycal \kappa(y,A)\,\mu(\dd y).
\end{equation}
Kernels also compose: $(\kappa_2\kappa_1)(y,A)=\int\kappa_2(x,A)\,
\kappa_1(y,\dd x)$ is the Chapman--Kolmogorov rule. This algebra is what
makes an autoregressive factorization
$p(x_{1:T})=\prod_t p(x_t\given x_{<t})$ or a Markov chain
$p(x_{0:T})=p_0\prod_t \kappa(x_{t-1},\dd x_t)$ a genuine probability
measure on the product space rather than a formal product of symbols
(Sections~\ref{sec:autoreg}, \ref{sec:markov-diffusion}).
\end{remark}

\section{Regular conditional distributions}

\begin{definition}[Regular conditional distribution]\label{def:rcd}
Let $X:\Omega\to\Xcal$ and $Y:\Omega\to\Ycal$ be random variables. A
\emph{regular conditional distribution} (r.c.d.) of $X$ given $Y$ is a
Markov kernel $\kappa$ from $\Ycal$ to $\Xcal$ such that for every
$A\in\Bcal_\Xcal$,
\begin{equation}\label{eq:rcd-defining}
  \kappa\bigl(Y,A\bigr)\;\eqas\;\cProb{X\in A}{Y}.
\end{equation}
Equivalently, for all $A\in\Bcal_\Xcal$, $B\in\Bcal_\Ycal$,
\begin{equation}\label{eq:rcd-joint}
  \Prob(X\in A,\,Y\in B)\;=\;\int_B \kappa(y,A)\,P_Y(\dd y).
\end{equation}
We then write $P_{X\given Y}(A\given y)\defeq\kappa(y,A)$.
\end{definition}

The point of~\eqref{eq:rcd-defining} is the order of quantifiers. For each
fixed $A$ the right-hand side is only an a.s.-defined random variable; the
r.c.d.\ asserts the existence of a \emph{single} kernel that
simultaneously serves as a version for \emph{all} $A$ at once, \emph{and}
is a genuine probability measure in $A$ for every $y$. That such a
reconciliation is possible is not automatic.

\begin{theorem}[Existence on standard Borel spaces]\label{thm:rcd-exists}
If $\Xcal$ is a standard Borel space (a Borel subset of a Polish space with
its Borel $\sigma$-algebra), then a regular conditional distribution of $X$
given $Y$ exists, and it is unique in the sense that any two r.c.d.s
$\kappa,\kappa'$ satisfy $\kappa(y,\cdot)=\kappa'(y,\cdot)$ for
$P_Y$-almost every $y$.
\end{theorem}

\begin{remark}[Why the topological hypothesis is unavoidable]
Without a hypothesis like ``standard Borel'' the r.c.d.\ can fail to exist:
one can patch together the a.s.-defined versions $\cProb{X\in A}{Y}$ over a
countable generating algebra, but extending to a measure for \emph{every}
$y$ requires controlling uncountably many null sets at once, which needs
the regularity that Polish spaces provide (via, e.g., the existence of
a countable basis and the Carath\'eodory extension along a compact class).
For the ML reader the practical upshot is reassuring: every space that
occurs in practice --- $\R^d$, compact manifolds, sequence spaces,
$\{0,1\}^n$, spaces of measures with the weak topology --- is standard
Borel, so the conditional distributions written in papers do exist. The
hypothesis is exactly the ``standing convention'' announced in the Preface.
\end{remark}

\begin{theorem}[Disintegration]\label{thm:disintegration}
Let $\mu$ be the joint law of $(X,Y)$ on $\Xcal\times\Ycal$ with $\Xcal$
standard Borel, and let $\nu=P_Y$ be the $Y$-marginal. Then there is a
$\nu$-a.e.-unique Markov kernel $\kappa(y,\dd x)$ with
\begin{equation}\label{eq:disint}
  \mu(\dd x,\dd y)\;=\;\kappa(y,\dd x)\,\nu(\dd y),
\end{equation}
meaning $\displaystyle\int f\,\dd\mu=\int_\Ycal\!\int_\Xcal
f(x,y)\,\kappa(y,\dd x)\,\nu(\dd y)$ for all bounded measurable $f$. The
kernel $\kappa$ is the r.c.d.\ of $X$ given $Y$.
\end{theorem}

Disintegration~\eqref{eq:disint} is the theorem that justifies writing a
joint density as ``conditional $\times$ marginal'' --- the single most
common factorization in machine learning --- at the level of measures,
before any density exists. With it, conditional expectation is recovered as
an honest integral against a measure:
\begin{equation}\label{eq:ce-via-rcd}
  \cE{X}{Y=y}\;=\;\int_\Xcal x\,\kappa(y,\dd x),
  \qquad
  \cE{\varphi(X)}{Y=y}\;=\;\int_\Xcal \varphi(x)\,\kappa(y,\dd x),
\end{equation}
valid for $P_Y$-a.e.\ $y$ whenever the integrals converge. This is the form
in which a mathematician should \emph{read} every ``$\E_{x\sim p(\cdot\given
y)}[\varphi(x)]$'' in a paper.

\chapter{Conditional densities and Bayes' rule}\label{ch:cond-density}

A density is a derivative of one measure with respect to another. ML uses
the letter $p$ for all of them, so the reference measure must be supplied
by the reader.

\section{When a conditional density exists}

\begin{definition}[Conditional density]\label{def:cond-density}
Fix a $\sigma$-finite reference measure $\lambda$ on $\Xcal$ (Lebesgue on
$\R^d$, counting on a discrete set). Suppose the joint law of $(X,Y)$ has a
density $p_{X,Y}$ with respect to $\lambda\otimes\rho$ for some
$\sigma$-finite $\rho$ on $\Ycal$, and set
$p_Y(y)=\int_\Xcal p_{X,Y}(x,y)\,\lambda(\dd x)$, the $Y$-marginal density
w.r.t.\ $\rho$. Then for $y$ with $p_Y(y)>0$ define
\begin{equation}\label{eq:cond-density}
  p_{X\given Y}(x\given y)\;\defeq\;\frac{p_{X,Y}(x,y)}{p_Y(y)} ,
\end{equation}
and (arbitrarily, e.g.\ as a fixed reference density) where $p_Y(y)=0$.
\end{definition}

\begin{proposition}[The density is a density of the r.c.d.]
\label{prop:density-is-rcd}
Under Definition~\ref{def:cond-density}, the kernel
$\kappa(y,A)=\int_A p_{X\given Y}(x\given y)\,\lambda(\dd x)$ is a regular
conditional distribution of $X$ given $Y$. In particular
$\{y:p_Y(y)=0\}$ is $P_Y$-null, so the arbitrary choice there is immaterial,
and~\eqref{eq:ce-via-rcd} becomes
$\cE{\varphi(X)}{Y=y}=\int\varphi(x)\,p_{X\given Y}(x\given y)\,\lambda(\dd x)$.
\end{proposition}
\begin{proof}
Verify~\eqref{eq:rcd-joint}: for $A\in\Bcal_\Xcal$, $B\in\Bcal_\Ycal$,
Tonelli gives
$\int_B\!\int_A p_{X\given Y}(x\given y)\lambda(\dd x)\,p_Y(y)\rho(\dd y)
=\int_{A\times B}p_{X,Y}\,\dd(\lambda\otimes\rho)=\Prob(X\in A,Y\in B)$,
using $p_{X\given Y}\,p_Y=p_{X,Y}$ off the null set $\{p_Y=0\}$.
\end{proof}

\begin{warning}[The same letter $p$ is three different functions]
In~\eqref{eq:cond-density} the symbols $p_{X,Y}$, $p_Y$, and
$p_{X\given Y}$ are distinct functions on distinct domains; papers drop the
subscripts and write $p(x,y)$, $p(y)$, $p(x\given y)$, relying on the
\emph{argument names} to disambiguate. Thus in ML ``$p(y)$'' and
``$p(x)$'' may denote different functions of a real variable. A
mathematician should mentally restore the subscripts: the bar in
$p(x\given y)$ encodes ``the $\Xcal$-density of the r.c.d.\ at parameter
$y$,'' not a quotient of values of one fixed function.
\end{warning}

\section{Bayes' rule as an identity of densities}

\begin{proposition}[Bayes' rule]\label{prop:bayes}
With the hypotheses of Definition~\ref{def:cond-density} and $p_X(x)>0$,
where $p_X(x)=\int p_{X,Y}(x,y)\,\rho(\dd y)$ is the $X$-marginal,
\begin{equation}\label{eq:bayes}
  p_{Y\given X}(y\given x)
   \;=\;\frac{p_{X\given Y}(x\given y)\,p_Y(y)}{p_X(x)}
   \;=\;\frac{p_{X\given Y}(x\given y)\,p_Y(y)}
            {\displaystyle\int_\Ycal p_{X\given Y}(x\given y')\,p_Y(y')\,\rho(\dd y')}.
\end{equation}
\end{proposition}

Read with $Y=\theta$ a parameter and $X=\Dcal$ the data, \eqref{eq:bayes}
is the posterior $=$ likelihood $\times$ prior $/$ evidence formula of
Bayesian ML (Section~\ref{sec:posterior}). Read with $Y$ a class label it
is the generative classifier. The denominator --- the \emph{evidence} or
\emph{partition function} --- is the normalizing integral whose
intractability is the reason variational and Monte-Carlo machinery exists
at all.

\section{The Borel--Kolmogorov paradox}\label{sec:bk-paradox}

We can now state precisely the ambiguity warned of in
Chapter~\ref{ch:why}, and its resolution.

\begin{example}[Conditioning on a great circle vs.\ a meridian]
\label{ex:bk}
Let $X$ be uniform on the unit sphere $S^2$. Consider the event ``$X$ lies
on a fixed great circle $C$.'' One may present $C$ as
$\{$longitude $=\,$const$\}$ (a meridian) or, after rotating, as
$\{$latitude $=0\}$ (the equator). Conditioning $X$ on the meridian, the
r.c.d.\ along $C$ is \emph{not} uniform in arc length (it is weighted by
$\cos(\text{latitude})$); conditioning on the equator, it \emph{is}
uniform. Both are correct r.c.d.s for their respective \emph{variables},
yet they assign different ``conditional laws'' to the same null set $C$.
\end{example}

\begin{remark}[Diagnosis and resolution]
There is no contradiction with Theorem~\ref{thm:rcd-exists}. An r.c.d.\ is a
kernel attached to a \emph{random variable} $Y$ (longitude, or latitude),
not to a bare null event. Uniqueness in Theorem~\ref{thm:rcd-exists} is
$P_Y$-a.e.\ \emph{for that $Y$}; changing $Y$ changes the kernel and there
is no reason two different conditioning variables should agree on a single
fibre that both happen to contain. The lesson for reading ML papers:
\begin{quote}
``$p(x\given y)$ at a particular $y$'' is meaningful only relative to the
\emph{family} $\{p(\cdot\given y)\}_y$ --- equivalently, relative to the
disintegrating variable $Y$. Conditioning on an isolated measure-zero
event, detached from such a family, is undefined.
\end{quote}
This is why diffusion and flow papers are careful to define a whole path of
conditional laws $p_t(\cdot\given y)$ indexed by $t$ and $y$, rather than
conditioning on individual zero-probability level sets ad hoc.
\end{remark}

% =====================================================================
\part{Conditional Independence}\label{part:cond-indep}
% =====================================================================

\chapter{Conditional independence}\label{ch:cond-indep}

The symbol $X\indep Y\given Z$ --- ``$X$ and $Y$ are independent given
$Z$'' --- is the load-bearing assumption of graphical models, causal
inference, and almost every probabilistic factorization in ML. It has
several equivalent definitions; knowing which one a paper is using removes
most confusion.

\section{Equivalent definitions}

\begin{definition}[Conditional independence]\label{def:ci}
Let $X,Y,Z$ be random variables into standard Borel spaces. We say $X$ and
$Y$ are \emph{conditionally independent given $Z$}, written
$X\indep Y\given Z$, if any one of the following equivalent conditions
holds.
\begin{enumerate}[label=\textup{(CI\arabic*)},leftmargin=3.4em]
\item \textup{(Kernel factorization)} a regular conditional distribution of
      $(X,Y)$ given $Z$ factors:
      \begin{equation}\label{eq:ci-kernel}
        P_{(X,Y)\given Z}(\dd x,\dd y\given z)
        \;=\;P_{X\given Z}(\dd x\given z)\,P_{Y\given Z}(\dd y\given z)
        \quad\text{for }P_Z\text{-a.e.\ }z.
      \end{equation}
\item \textup{(Conditional-probability product)} for all
      $A\in\Bcal_\Xcal$, $B\in\Bcal_\Ycal$,
      \begin{equation}\label{eq:ci-prob}
        \cProb{X\in A,\,Y\in B}{Z}\;\eqas\;
        \cProb{X\in A}{Z}\,\cProb{Y\in B}{Z}.
      \end{equation}
\item \textup{(Conditional-expectation product)} for all bounded Borel
      $f,g$,
      \begin{equation}\label{eq:ci-exp}
        \cE{f(X)g(Y)}{Z}\;\eqas\;\cE{f(X)}{Z}\,\cE{g(Y)}{Z}.
      \end{equation}
\item \textup{(Sufficiency form)} for all bounded Borel $f$,
      \begin{equation}\label{eq:ci-suff}
        \cE{f(X)}{Y,Z}\;\eqas\;\cE{f(X)}{Z};
      \end{equation}
      i.e.\ once $Z$ is known, $Y$ carries no further information about $X$.
\item \textup{(Density factorization, when densities exist)}
      \begin{equation}\label{eq:ci-density}
        p_{X,Y\given Z}(x,y\given z)\;=\;p_{X\given Z}(x\given z)\,
        p_{Y\given Z}(y\given z)\qquad(\text{a.e.}).
      \end{equation}
\end{enumerate}
\end{definition}

\begin{proposition}[Equivalence]
Conditions \textup{(CI1)--(CI5)} are equivalent (with \textup{(CI5)}
requiring the existence of the relevant densities).
\end{proposition}
\begin{proof}[Notes]
(CI1)$\Leftrightarrow$(CI2) is~\eqref{eq:rcd-defining} read for the pair
$(X,Y)$ and the product kernel. (CI2)$\Leftrightarrow$(CI3) extends
indicators to bounded functions by linearity and monotone convergence.
(CI3)$\Rightarrow$(CI4): for bounded $f$ and any bounded $g(Y),h(Z)$,
$\E[\cE{f(X)}{Y,Z}g(Y)h(Z)]=\E[f(X)g(Y)h(Z)]
=\E[h(Z)\cE{f(X)g(Y)}{Z}] =\E[h(Z)\cE{f(X)}{Z}\cE{g(Y)}{Z}]
=\E[\cE{f(X)}{Z}g(Y)h(Z)]$, using (CI3) and taking out $\cE{g(Y)}{Z}$;
since products $g(Y)h(Z)$ generate $\sigma(Y,Z)$ this
gives~\eqref{eq:ci-suff}. (CI4)$\Rightarrow$(CI3) reverses the computation
via the tower rule. (CI5) is~\eqref{eq:ci-kernel} written through
densities (Proposition~\ref{prop:density-is-rcd}).
\end{proof}

\begin{remark}[The symmetry that the bar hides]
Form~\eqref{eq:ci-suff} looks asymmetric in $X$ and $Y$, but conditional
independence is symmetric: $X\indep Y\given Z\iff Y\indep X\given Z$. The
sufficiency reading is nonetheless the one most used in ML --- ``$Z$ is a
sufficient statistic,'' ``the representation $Z$ makes the label $Y$
conditionally independent of the input $X$,'' the Markov property of latent
states --- so it is worth recognizing it as merely one face of
Definition~\ref{def:ci}.
\end{remark}

\section{The graphoid axioms}

Conditional independence behaves like a calculus of ``irrelevance.'' The
following identities, the \emph{graphoid axioms}, are the rules that
graphical-model algorithms (d-separation, message passing) silently rely
on. Write $X\indep Y\given Z$ for the ternary relation.

\begin{proposition}[Semi-graphoid axioms]\label{prop:graphoid}
For random variables $W,X,Y,Z$:
\begin{align}
&\textup{(Symmetry)} && X\indep Y\given Z \;\Longrightarrow\; Y\indep X\given Z;\\
&\textup{(Decomposition)} && X\indep (Y,W)\given Z \;\Longrightarrow\; X\indep Y\given Z;\\
&\textup{(Weak union)} && X\indep (Y,W)\given Z \;\Longrightarrow\; X\indep Y\given (Z,W);\\
&\textup{(Contraction)} && \bigl(X\indep Y\given Z\bigr)\,\wedge\,
   \bigl(X\indep W\given (Y,Z)\bigr)\;\Longrightarrow\; X\indep (Y,W)\given Z.
\end{align}
\end{proposition}

\begin{proposition}[Intersection, under positivity]\label{prop:intersection}
If the joint density is strictly positive (more generally, under a
suitable positivity condition), the additional \emph{intersection} axiom
holds:
\begin{equation}
  \bigl(X\indep Y\given (Z,W)\bigr)\,\wedge\,
  \bigl(X\indep W\given (Z,Y)\bigr)\;\Longrightarrow\;
  X\indep (Y,W)\given Z.
\end{equation}
\end{proposition}

\begin{warning}[Intersection can fail]
The intersection axiom is the one that requires a hypothesis. Without
positivity it is false: take $Y=W$ a single fair coin and $X=Y$; then
$X\indep Y\given (Z,W)$ and $X\indep W\given (Z,Y)$ hold trivially (the
conditioning already determines $X$), yet $X$ is not independent of
$(Y,W)$ given $Z$. ML papers that invoke ``by the graphoid axioms'' while
manipulating sparse or deterministic models are implicitly assuming the
positivity that licenses intersection; a careful reader checks it.
\end{warning}

\section{Factorization over graphs}

\begin{definition}[Bayesian-network factorization]\label{def:bn}
Let $G$ be a directed acyclic graph on vertices $\{1,\dots,d\}$ with parent
sets $\mathrm{pa}(i)$. A law $p$ on $\Xcal_1\times\cdots\times\Xcal_d$
\emph{factorizes according to $G$} if
\begin{equation}\label{eq:bn-factor}
  p(x_1,\dots,x_d)\;=\;\prod_{i=1}^d p\bigl(x_i\given x_{\mathrm{pa}(i)}\bigr).
\end{equation}
\end{definition}

\begin{theorem}[Factorization $\Leftrightarrow$ Markov property]
\label{thm:bn-markov}
A positive law factorizes as~\eqref{eq:bn-factor} if and only if it
satisfies the local Markov property of $G$: each variable is conditionally
independent of its non-descendants given its parents,
$X_i\indep X_{\mathrm{nd}(i)\setminus\mathrm{pa}(i)}\given X_{\mathrm{pa}(i)}$.
Either condition is equivalent to the global Markov property
(d-separation $\Rightarrow$ conditional independence).
\end{theorem}

Equation~\eqref{eq:bn-factor} is the rigorous meaning of the ``plate
diagrams'' and factor graphs that open most probabilistic-ML papers, and
each factor $p(x_i\given x_{\mathrm{pa}(i)})$ is a Markov kernel in the
sense of Definition~\ref{def:kernel}. The autoregressive models of
Section~\ref{sec:autoreg} are the special case of the complete DAG with
$\mathrm{pa}(i)=\{1,\dots,i-1\}$, where~\eqref{eq:bn-factor} is not an
assumption at all but the chain rule of probability.

% =====================================================================
\part{Conditional Information Measures}\label{part:info}
% =====================================================================

\chapter{Relative entropy and the conditional information measures}
\label{ch:info}

Every information-theoretic quantity in ML is built from one primitive ---
relative entropy --- so we start there and obtain entropy, cross-entropy,
mutual information, and their conditional versions as corollaries. Logarithms
are natural unless a paper says ``bits.''

\section{Relative entropy (KL divergence)}

\begin{definition}[Kullback--Leibler divergence]\label{def:kl}
For probability measures $P,Q$ on a common $(\Xcal,\Bcal)$,
\begin{equation}\label{eq:kl}
  \KL(P\,\|\,Q)\;\defeq\;
  \begin{cases}
    \displaystyle\int_\Xcal \log\frac{\dd P}{\dd Q}\,\dd P
      \;=\;\int_\Xcal \frac{\dd P}{\dd Q}\log\frac{\dd P}{\dd Q}\,\dd Q,
      & P\ll Q,\\[1.2em]
    +\infty, & P\not\ll Q.
  \end{cases}
\end{equation}
If $P,Q$ have densities $p,q$ w.r.t.\ a common reference $\lambda$, then
$\KL(P\|Q)=\int p\log(p/q)\,\dd\lambda$.
\end{definition}

\begin{proposition}[Gibbs' inequality]\label{prop:gibbs}
$\KL(P\|Q)\ge 0$, with equality iff $P=Q$. Consequently $\KL$ is a
\emph{divergence}: it separates points. It is \emph{not} a metric --- it is
asymmetric and violates the triangle inequality.
\end{proposition}
\begin{proof}
Apply Jensen to the strictly convex $-\log$:
\[
  -\KL(P\|Q)=\int\log\frac{\dd Q}{\dd P}\,\dd P
   \;\le\;\log\!\int\frac{\dd Q}{\dd P}\,\dd P=\log Q(\Xcal)=0,
\]
with equality iff $\dd Q/\dd P$ is $P$-a.s.\ constant, i.e.\ $P=Q$.
\end{proof}

\begin{warning}[The absolute-continuity trap]
$\KL(P\|Q)=+\infty$ whenever $P$ places mass where $Q$ does not. This is
the rigorous reason a generative model with a too-small support, or a
variational posterior $q$ that is more spread out than the true posterior,
incurs the notorious ``mode covering vs.\ mode seeking'' asymmetry:
$\KL(q\|p)$ (used in variational inference, Section~\ref{sec:elbo}) blows
up if $q$ has mass where $p\approx0$, whereas $\KL(p\|q)$ (used in maximum
likelihood, Section~\ref{sec:mle}) blows up if $q$ misses mass that $p$
has. Same symbol $\KL$, opposite failure modes, decided entirely by
argument order.
\end{warning}

\section{Entropy and cross-entropy}

\begin{definition}[Entropy and cross-entropy]\label{def:entropy}
For a discrete law $P$ with pmf $p$ on a countable $\Xcal$,
\begin{equation}
  \Ent(P)\;=\;\Ent(X)\;\defeq\;-\!\sum_{x}p(x)\log p(x)
            \;=\;-\,\E_{P}[\log p(X)],
\end{equation}
and for $P\ll\lambda$ with density $p$ (e.g.\ $\lambda=$ Lebesgue), the
\emph{differential entropy}
\begin{equation}
  \diffent(X)\;\defeq\;-\!\int_\Xcal p(x)\log p(x)\,\lambda(\dd x).
\end{equation}
The \emph{cross-entropy} of $Q$ (density $q$) relative to $P$ is
\(
  \Ent(P,Q)\defeq-\E_P[\log q(X)].
\)
\end{definition}

\begin{proposition}[The bookkeeping identity]\label{prop:ce-kl}
Whenever the terms are finite,
\begin{equation}\label{eq:ce-decomp}
  \Ent(P,Q)\;=\;\Ent(P)\;+\;\KL(P\,\|\,Q).
\end{equation}
\end{proposition}

Equation~\eqref{eq:ce-decomp} is why ``minimizing cross-entropy'' and
``minimizing KL to the data distribution'' are the same training objective:
$\Ent(P)$ does not depend on the model $Q$, so
$\argmin_Q \Ent(P,Q)=\argmin_Q\KL(P\|Q)$. We return to this in
Section~\ref{sec:mle}.

\begin{warning}[Differential entropy is not entropy]\label{warn:diffent}
For continuous $X$, $\diffent(X)$ can be negative, is \emph{not} invariant
under smooth reparameterization (under $X\mapsto T(X)$ it changes by
$\E[\log|\det\nabla T|]$), and is not the limit of discrete entropies of
finer quantizations (that limit diverges). It is best regarded only as a
term that makes other quantities --- mutual information, KL --- behave well,
since those are differences in which the offending Jacobians and infinities
cancel. A mathematician should distrust any standalone use of $\diffent$.
\end{warning}

\section{Conditional entropy}

\begin{definition}[Conditional entropy]\label{def:cond-entropy}
For discrete $(X,Y)$,
\begin{equation}\label{eq:cond-entropy}
  \Ent(X\given Y)\;\defeq\;\sum_{y}p_Y(y)\,\Ent(X\given Y=y)
   \;=\;-\!\sum_{x,y}p_{X,Y}(x,y)\,\log p_{X\given Y}(x\given y),
\end{equation}
where $\Ent(X\given Y=y)$ is the entropy of the fibre law
$P_{X\given Y}(\cdot\given y)$. Equivalently,
\[
  \Ent(X\given Y)=\E_Y\bigl[\Ent\bigl(P_{X\given Y}(\cdot\given Y)\bigr)\bigr].
\]
\end{definition}

So conditional entropy is the \emph{expected entropy of the regular
conditional distribution} --- an average over $Y$ of an entropy computed
inside each fibre, exactly the disintegration~\eqref{eq:disint} at work.
It is \emph{not} the entropy of a conditional law at a single $y$; the
single-$y$ object is $\Ent(X\given Y=y)$, and the unadorned $\Ent(X\given
Y)$ has already integrated it out.

\begin{proposition}[Chain rule and conditioning reduces entropy]
\label{prop:entropy-chain}
For discrete variables,
\begin{equation}\label{eq:entropy-chain}
  \Ent(X,Y)\;=\;\Ent(Y)+\Ent(X\given Y),\qquad
  0\;\le\;\Ent(X\given Y)\;\le\;\Ent(X),
\end{equation}
the right inequality (``conditioning cannot increase entropy,'' with
equality iff $X\indep Y$) holding \emph{on average} only --- it can happen
that $\Ent(X\given Y=y)>\Ent(X)$ for particular $y$.
\end{proposition}

\section{Mutual information and its conditional form}

\begin{definition}[Mutual and conditional mutual information]\label{def:mi}
\begin{align}
  \MI(X;Y)&\;\defeq\;\KL\bigl(P_{X,Y}\,\big\|\,P_X\otimes P_Y\bigr),\label{eq:mi}\\
  \MI(X;Y\given Z)&\;\defeq\;
    \E_Z\Bigl[\KL\bigl(P_{X,Y\given Z}(\cdot\given Z)\,\big\|\,
       P_{X\given Z}(\cdot\given Z)\otimes P_{Y\given Z}(\cdot\given Z)\bigr)\Bigr].
  \label{eq:cmi}
\end{align}
\end{definition}

Definition~\eqref{eq:mi} is reference-measure-free and so escapes the
pathologies of differential entropy (Warning~\ref{warn:diffent}): mutual
information is a KL divergence between honest measures and is therefore
nonnegative, invariant under bijective reparameterization of $X$ and of
$Y$, and equal to $+\infty$ exactly when $P_{X,Y}\not\ll P_X\otimes P_Y$.
The conditional version~\eqref{eq:cmi} is, once more, an \emph{average over
$Z$} of a per-fibre KL.

\begin{proposition}[Identities]\label{prop:mi-ids}
When the entropies are finite,
\begin{align}
  \MI(X;Y)&=\Ent(X)-\Ent(X\given Y)=\Ent(Y)-\Ent(Y\given X)\ \ge 0,\\
  \MI(X;Y\given Z)&=\Ent(X\given Z)-\Ent(X\given Y,Z)\ \ge 0,\\
  \MI(X;Y,Z)&=\MI(X;Z)+\MI(X;Y\given Z) \quad\text{(chain rule).}
  \label{eq:mi-chain}
\end{align}
Moreover $\MI(X;Y\given Z)=0\iff X\indep Y\given Z$, tying
Part~\ref{part:cond-indep} to Part~\ref{part:info}: \emph{conditional
independence is exactly the vanishing of conditional mutual information.}
\end{proposition}

\begin{remark}[Where these appear]
Conditional mutual information $\MI(X;Y\given Z)$ is the quantity behind the
information-bottleneck objective, the ``minimality/sufficiency'' analyses of
learned representations, and the conditional-independence tests of causal
discovery. The chain rule~\eqref{eq:mi-chain} is the information-theoretic
shadow of the probability chain rule that drives autoregressive models
(Section~\ref{sec:autoreg}). Recognizing all of these as averaged
per-fibre KL divergences --- never as arithmetic on single-$y$ entropies ---
is the main thing a mathematician needs to read them safely.
\end{remark}

\section{Conditional KL and the chain rule for divergence}

\begin{definition}[Conditional / expected KL]\label{def:cond-kl}
Given kernels $P_{X\given Y},Q_{X\given Y}$ and a law $P_Y$,
\begin{equation}\label{eq:cond-kl}
  \KL\bigl(P_{X\given Y}\,\big\|\,Q_{X\given Y}\,\big|\,P_Y\bigr)
   \;\defeq\;\int_\Ycal
     \KL\bigl(P_{X\given Y}(\cdot\given y)\,\big\|\,Q_{X\given Y}(\cdot\given y)\bigr)
     \,P_Y(\dd y).
\end{equation}
\end{definition}

\begin{proposition}[Chain rule for KL]\label{prop:kl-chain}
For joint laws $P=P_Y\,P_{X\given Y}$ and $Q=Q_Y\,Q_{X\given Y}$,
\begin{equation}\label{eq:kl-chain}
  \KL(P_{X,Y}\,\|\,Q_{X,Y})
   \;=\;\KL(P_Y\,\|\,Q_Y)
       \;+\;\KL\bigl(P_{X\given Y}\,\big\|\,Q_{X\given Y}\,\big|\,P_Y\bigr).
\end{equation}
\end{proposition}

The conditional KL~\eqref{eq:cond-kl} --- note it averages the second
argument's null-set behavior over $P_Y$, the \emph{first} law's marginal ---
is the per-step training signal of diffusion models, where the loss is a
sum of terms $\KL\bigl(q(x_{t-1}\given x_t,x_0)\,\|\,p_\theta(x_{t-1}\given
x_t)\bigr)$ averaged over the forward process
(Section~\ref{sec:markov-diffusion}). The chain rule~\eqref{eq:kl-chain} is
what telescopes those per-step divergences into a bound on the divergence
between full trajectory laws.

% =====================================================================
\part{The Machine-Learning Symbol Dictionary}\label{part:dictionary}
% =====================================================================

This part is the reverse map. Each entry takes a symbol \emph{as it is
actually printed in a thesis}, states what a reader should mentally
substitute, and points to the rigorous object in
Parts~\ref{part:foundations}--\ref{part:info}. The recurring move is
always one of three: a bar denotes a Markov kernel
(Definition~\ref{def:kernel}), a conditional expectation
(Definition~\ref{def:cond-exp}), or a conditional density
(Definition~\ref{def:cond-density}); everything else is bookkeeping.

\chapter{Likelihood, posterior, and the variational bar}

\section{Likelihood, NLL, and MLE as a KL projection}\label{sec:mle}

\begin{mlentry}{$p_\theta(x)$,\quad $p_\theta(y\given x)$,\quad
$\mathcal L(\theta)=\prod_i p_\theta(x_i)$,\quad
$-\log p_\theta(y\given x)$.}
A \emph{model} is a parameterized family of probability measures
$\{P_\theta:\theta\in\Theta\}$ on the data space, each with density
$p_\theta$ w.r.t.\ a fixed reference $\lambda$. Read $p_\theta(x)$ as ``the
density of $P_\theta$ at $x$,'' a function of $x$ for fixed $\theta$; read
the \emph{likelihood} $\theta\mapsto p_\theta(x)$ as the same expression
viewed as a function of $\theta$ for fixed observed $x$. The
\emph{conditional likelihood} $p_\theta(y\given x)$ is a Markov kernel
(Definition~\ref{def:kernel}) indexed by $\theta$: a discriminative model
predicting $y$ from $x$. The \emph{negative log-likelihood} (NLL)
$-\log p_\theta(y\given x)$ is the per-example loss of supervised learning.
\end{mlentry}

\begin{proposition}[Maximum likelihood is KL projection onto the model]
\label{prop:mle-kl}
Let $\widehat P_n=\frac1n\sum_{i=1}^n\delta_{x_i}$ be the empirical
distribution of i.i.d.\ data $x_i\sim P_\star$. Then
\begin{equation}
  \argmax_{\theta}\ \frac1n\sum_{i=1}^n\log p_\theta(x_i)
  \;=\;\argmin_{\theta}\ \Ent\!\bigl(\widehat P_n,\,P_\theta\bigr)
  \;=\;\argmin_{\theta}\ \KL\!\bigl(\widehat P_n\,\big\|\,P_\theta\bigr),
\end{equation}
and as $n\to\infty$ the population objective is
$\argmin_\theta \KL(P_\star\|P_\theta)$.
\end{proposition}
\begin{proof}
Maximizing the average log-likelihood equals minimizing the cross-entropy
$\Ent(\widehat P_n,P_\theta)=-\frac1n\sum_i\log p_\theta(x_i)$; subtract the
$\theta$-free term $\Ent(\widehat P_n)$ and apply~\eqref{eq:ce-decomp}.
\end{proof}

So ``train by minimizing cross-entropy'' \emph{is} ``project the data
distribution onto the model in KL'' --- the forward $\KL(P_\star\|P_\theta)$,
the mode-covering one of Warning~\ref{warn:diffent}'s companion. The bar in
$p_\theta(y\given x)$ contributes the conditional version: minimizing
$\E_{x}\bigl[\KL\bigl(P_\star(\cdot\given x)\,\|\,P_\theta(\cdot\given
x)\bigr)\bigr]$, the conditional KL~\eqref{eq:cond-kl} of
Definition~\ref{def:cond-kl}.

\section{Posterior and evidence}\label{sec:posterior}

\begin{mlentry}{$p(\theta\given\Dcal)$,\quad $p(\Dcal\given\theta)$,\quad
$p(\Dcal)=\int p(\Dcal\given\theta)p(\theta)\dd\theta$.}
In the Bayesian picture the parameter $\theta$ is itself a random variable
with prior law $p(\theta)$. The \emph{posterior} $p(\theta\given\Dcal)$ is
literally the conditional density of Definition~\ref{def:cond-density} with
$X=\theta$, $Y=\Dcal$, supplied by Bayes' rule
(Proposition~\ref{prop:bayes}):
\begin{equation}\label{eq:posterior}
  p(\theta\given\Dcal)\;=\;\frac{p(\Dcal\given\theta)\,p(\theta)}{p(\Dcal)},
  \qquad p(\Dcal)=\int_\Theta p(\Dcal\given\theta)\,p(\theta)\,\dd\theta.
\end{equation}
The normalizer $p(\Dcal)$ --- the \emph{evidence} or \emph{marginal
likelihood} --- is the $\Dcal$-marginal density and is the integral whose
intractability motivates everything in the next section.
\end{mlentry}

\begin{warning}[$p(x;\theta)$ vs.\ $p(x\given\theta)$: semicolon is not a bar]
A frequentist writes $p(x;\theta)$ or $p_\theta(x)$ --- $\theta$ is a fixed
index, not an event, and there is no measure on $\Theta$. A Bayesian writes
$p(x\given\theta)$ --- $\theta$ is a random variable and the bar is genuine
conditioning. The two coincide \emph{numerically} but mean different things:
only the second makes $p(\theta\given x)$ meaningful. When a paper slides
between $p(x;\theta)$ and $p(x\given\theta)$ it is silently choosing whether
$\Theta$ carries a probability measure. A mathematician should track which.
\end{warning}

\section{The variational posterior and the ELBO}\label{sec:elbo}

\begin{mlentry}{$q_\phi(z\given x)$,\quad
$\E_{q_\phi(z\given x)}[\cdot]$,\quad ELBO.}
A latent-variable model posits $p_\theta(x,z)=p_\theta(x\given z)\,p(z)$
with a prior $p(z)$ and a decoder kernel $p_\theta(x\given z)$. The exact
posterior $p_\theta(z\given x)=p_\theta(x,z)/p_\theta(x)$ is intractable
because $p_\theta(x)=\int p_\theta(x\given z)p(z)\dd z$ is. The
\emph{variational posterior} (or encoder, or recognition kernel)
$q_\phi(z\given x)$ is a second Markov kernel, parameterized by $\phi$,
introduced to approximate it. The expectation $\E_{q_\phi(z\given x)}[\cdot]$
is the integral $\int(\cdot)\,q_\phi(\dd z\given x)$ against that kernel ---
a conditional expectation under $q$, in the sense of~\eqref{eq:ce-via-rcd}.
\end{mlentry}

\begin{proposition}[Evidence lower bound and its gap]\label{prop:elbo}
For any kernel $q_\phi(\cdot\given x)\ll p_\theta(\cdot\given x)$,
\begin{equation}\label{eq:elbo}
  \log p_\theta(x)
  \;=\;\underbrace{\E_{q_\phi(z\given x)}\!\Bigl[\log\frac{p_\theta(x,z)}
        {q_\phi(z\given x)}\Bigr]}_{\displaystyle \mathrm{ELBO}(\theta,\phi;x)}
   \;+\;\KL\!\bigl(q_\phi(\cdot\given x)\,\big\|\,p_\theta(\cdot\given x)\bigr),
\end{equation}
and since the KL term is $\ge 0$, $\mathrm{ELBO}(\theta,\phi;x)\le\log
p_\theta(x)$ with equality iff $q_\phi(\cdot\given x)=p_\theta(\cdot\given
x)$. Equivalently
\begin{equation}\label{eq:elbo-split}
  \mathrm{ELBO}(\theta,\phi;x)
  =\E_{q_\phi(z\given x)}\!\bigl[\log p_\theta(x\given z)\bigr]
   -\KL\!\bigl(q_\phi(\cdot\given x)\,\big\|\,p(\cdot)\bigr).
\end{equation}
\end{proposition}
\begin{proof}
Insert $1=q_\phi(z\given x)/q_\phi(z\given x)$ into $\log
p_\theta(x)=\log\int p_\theta(x,z)\dd z$ is not valid directly; instead
write $\log p_\theta(x)=\E_{q_\phi(z\given x)}[\log p_\theta(x)]$ (the
integrand is constant in $z$), expand $\log p_\theta(x)=\log
p_\theta(x,z)-\log p_\theta(z\given x)$, and add and subtract $\log
q_\phi(z\given x)$. The two resulting expectations are the ELBO and the KL.
Equation~\eqref{eq:elbo-split} regroups $p_\theta(x,z)=p_\theta(x\given
z)p(z)$.
\end{proof}

\begin{remark}[Three bars, three different objects, one equation]
Equation~\eqref{eq:elbo} contains the encoder kernel $q_\phi(z\given x)$, the
decoder kernel $p_\theta(x\given z)$, and the exact-posterior kernel
$p_\theta(z\given x)$ --- three distinct Markov kernels, written with the
same bar. The variational gap is the \emph{reverse} KL
$\KL(q_\phi\|p_\theta(\cdot\given x))$, the mode-seeking direction; this is
why VAEs are prone to posterior collapse and under-dispersed encoders. The
mathematician's reading of ``maximize the ELBO'' is: minimize, over $\phi$,
the reverse conditional KL between two kernels, simultaneously fitting
$\theta$ by a lower bound on the marginal log-likelihood.
\end{remark}

\chapter{Sequential, diffusion, and reinforcement bars}

\section{Autoregressive models: the chain rule is not an assumption}
\label{sec:autoreg}

\begin{mlentry}{$p(x_{1:T})=\prod_{t=1}^T p(x_t\given x_{<t})$.}
For any random vector $(X_1,\dots,X_T)$ into standard Borel spaces, the
\emph{chain rule of probability} gives regular conditional distributions
$P_{X_t\given X_{<t}}$ with
\begin{equation}\label{eq:chain-prob}
  P_{X_{1:T}}(\dd x_{1:T})
   \;=\;\prod_{t=1}^T P_{X_t\given X_{<t}}(\dd x_t\given x_{<t}),
\end{equation}
by iterating disintegration (Theorem~\ref{thm:disintegration}) $T-1$ times.
Unlike a graphical-model factorization, \eqref{eq:chain-prob} holds for
\emph{every} joint law: it is an identity, not a modeling assumption. An
``autoregressive model'' replaces each true kernel by a learned kernel
$p_\theta(x_t\given x_{<t})$ (the softmax over a vocabulary in a language
model, a mixture or a logistic in a density model), and its NLL training
loss is the conditional cross-entropy summed over positions.
\end{mlentry}

\begin{remark}
What the model assumes is never the factorization~\eqref{eq:chain-prob} but
the \emph{parametric form and conditioning set} of each kernel: a finite
context window, a particular network, weight sharing across $t$. The bar in
$p(x_t\given x_{<t})$ is doing the honest work of
Theorem~\ref{thm:disintegration}; the modeling lives in the function class,
not in the conditioning.
\end{remark}

\section{Markov chains, transition kernels, and diffusion}
\label{sec:markov-diffusion}

\begin{mlentry}{$q(x_t\given x_{t-1})$,\quad $p_\theta(x_{t-1}\given x_t)$,\quad
$\nabla_x\log p_t(x)$,\quad $\nabla_x\log p_t(x\given y)$.}
A (time-inhomogeneous) Markov chain is a prior $p_0$ together with
transition kernels $\kappa_t(x_{t-1},\dd x_t)$; its law is the kernel
product of Remark~\ref{rem:kernel-compose}, and the Markov property is the
conditional independence $X_t\indep X_{<t-1}\given X_{t-1}$ of
Definition~\ref{def:ci}. A \emph{diffusion model} fixes a forward chain
$q(x_t\given x_{t-1})$ that gradually adds noise and learns a reverse chain
$p_\theta(x_{t-1}\given x_t)$; the training loss is a sum of conditional
KL terms~\eqref{eq:cond-kl}
$\KL\bigl(q(x_{t-1}\given x_t,x_0)\,\|\,p_\theta(x_{t-1}\given x_t)\bigr)$
averaged over the forward law, telescoped via the chain rule for
KL~\eqref{eq:kl-chain} into a bound on $-\log p_\theta(x_0)$.
\end{mlentry}

\begin{definition}[Score]
The \emph{(Stein) score} of a density $p_t$ on $\R^d$ is the gradient of its
log-density in the \emph{state}, $s_t(x)\defeq\nabla_x\log p_t(x)$ ---
\emph{not} a derivative in any parameter. It is invariant to the normalizing
constant of $p_t$, which is why score-based models can sidestep intractable
partition functions.
\end{definition}

\begin{proposition}[Conditional score and Bayes; the two guidances]
\label{prop:guidance}
For a joint $p_t(x,y)=p_t(x)\,p_t(y\given x)$ on $\R^d\times\Ycal$ with $y$
held fixed, differentiating $\log p_t(x\given y)=\log p_t(x)+\log
p_t(y\given x)-\log p_t(y)$ in $x$ gives the \emph{conditional score
decomposition}
\begin{equation}\label{eq:cond-score}
  \nabla_x\log p_t(x\given y)
  \;=\;\underbrace{\nabla_x\log p_t(x)}_{\text{unconditional score}}
   \;+\;\underbrace{\nabla_x\log p_t(y\given x)}_{\text{classifier gradient}} ,
\end{equation}
since $\nabla_x\log p_t(y)=0$. \emph{Classifier guidance} implements
\eqref{eq:cond-score} literally, adding the gradient of an auxiliary
classifier $p_t(y\given x)$. \emph{Classifier-free guidance} instead trains
one network to output both the conditional and unconditional scores
$s_\theta(x,y),s_\theta(x,\varnothing)$ and forms the guided field
\begin{equation}\label{eq:cfg}
  \widetilde s_\theta(x,y)
   \;=\;(1+w)\,s_\theta(x,y)\;-\;w\,s_\theta(x,\varnothing)
   \;=\;s_\theta(x,y)+w\bigl(s_\theta(x,y)-s_\theta(x,\varnothing)\bigr),
\end{equation}
with guidance weight $w\ge0$; the bracket is exactly the classifier
gradient $\nabla_x\log p_t(y\given x)$ of~\eqref{eq:cond-score}, so
$w=0$ recovers the true conditional score and $w>0$ samples from a
sharpened, tempered conditional $\propto p_t(x)\,p_t(y\given x)^{1+w}$.
\end{proposition}

\begin{remark}[A bar that is a gradient of a log of a kernel]
The symbol $\nabla_x\log p_t(x\given y)$ stacks three of this handbook's
constructions: a conditional density $p_t(x\given y)$
(Definition~\ref{def:cond-density}) indexed by a diffusion time $t$, its
logarithm, and a gradient in the \emph{conditioned} variable $x$.
Equation~\eqref{eq:cond-score} is just Bayes' rule
(Proposition~\ref{prop:bayes}) after a logarithm and a gradient --- the
single most important ``conditional thing'' in modern generative modeling,
and entirely elementary once each piece is named.
\end{remark}

\section{Policies, value functions, and the Bellman bar}\label{sec:rl}

\begin{mlentry}{$\pi(a\given s)$,\quad $V^\pi(s)=\E_\pi[\sum_t\gamma^t r_t\given S_0=s]$,\quad
$Q^\pi(s,a)$.}
A \emph{policy} $\pi(a\given s)$ is a Markov kernel from states to actions
(Definition~\ref{def:kernel}); together with the environment's transition
kernel $P(s'\given s,a)$ and reward it defines a law on trajectories. The
\emph{value function}
\begin{equation}\label{eq:value}
  V^\pi(s)\;\defeq\;\E_\pi\Bigl[\textstyle\sum_{t\ge0}\gamma^t R_t
     \,\Big|\,S_0=s\Bigr],\qquad
  Q^\pi(s,a)\;\defeq\;\E_\pi\Bigl[\textstyle\sum_{t\ge0}\gamma^t R_t
     \,\Big|\,S_0=s,A_0=a\Bigr],
\end{equation}
is an evaluated conditional expectation (Definition~\ref{def:eval-ce}) of a
discounted total reward; ``$\given S_0=s$'' is the conditioning-on-a-
random-variable of Chapter~\ref{ch:cond-on-rv}, well defined for $P_{S_0}$-a.e.\ $s$.
\end{mlentry}

\begin{proposition}[Bellman equation as the tower rule]\label{prop:bellman}
With the per-step conditional expectation over $A_0\sim\pi(\cdot\given s)$,
$S_1\sim P(\cdot\given s,A_0)$,
\begin{equation}\label{eq:bellman}
  V^\pi(s)\;=\;\E_{\pi}\bigl[\,R_0+\gamma\,V^\pi(S_1)\,\big|\,S_0=s\,\bigr].
\end{equation}
\end{proposition}
\begin{proof}
Split the geometric sum in~\eqref{eq:value} as
$R_0+\gamma\sum_{t\ge0}\gamma^t R_{t+1}$, and apply the tower
rule~\eqref{eq:tower} conditioning the tail on $(S_1,A_1,\dots)$ through
$S_1$; the Markov property collapses the tail's conditional expectation to
$V^\pi(S_1)$.
\end{proof}

So the Bellman equation --- the fixed-point identity behind value
iteration, $Q$-learning, and temporal-difference methods --- is nothing but
the tower property of conditional expectation
(Proposition~\ref{prop:ce-props}(4)) applied to a Markov reward process.
Every ``bootstrapping'' target $r+\gamma V(s')$ is a one-step conditional
expectation standing in for the full return.

\section{Filtrations and the martingale bar}

\begin{mlentry}{$\Fcal_t$,\quad $\E[X_{t+1}\given\Fcal_t]$,\quad
$\E[X_{t+1}\given\Fcal_t]=X_t$.}
A \emph{filtration} $(\Fcal_t)_{t\ge0}$ is an increasing family of
sub-$\sigma$-algebras, $\Fcal_s\subseteq\Fcal_t\subseteq\Fcal$ for $s\le
t$: ``the information available by time $t$.'' This is conditioning on a
$\sigma$-algebra (Definition~\ref{def:cond-exp}) in its purest form, with
no random variable to disintegrate against. A process is a
\emph{martingale} if it is adapted, integrable, and
$\E[X_{t+1}\given\Fcal_t]\eqas X_t$. Stochastic-approximation and SGD
convergence proofs are written in this language: the noise sequence is a
martingale-difference, $\E[\xi_{t+1}\given\Fcal_t]=0$, and the tower rule
plus a supermartingale convergence theorem do the rest.
\end{mlentry}

\begin{remark}[Why $\Fcal_t$ and not ``$\given\text{the past}$'']
Writing $\E[X_{t+1}\given\Fcal_t]$ rather than $\E[X_{t+1}\given
X_0,\dots,X_t]$ is the deliberately more general move: $\Fcal_t$ may encode
side information beyond the past values of $X$, and the conditional
expectation is defined directly on the $\sigma$-algebra without ever
choosing a conditioning random variable. For a mathematician this is the
\emph{cleanest} conditional bar in all of ML --- it is exactly
Definition~\ref{def:cond-exp}, untranslated.
\end{remark}

\chapter{A field guide to notational pitfalls}\label{ch:pitfalls}

We close the dictionary with the ambiguities that most often trip a
mathematician, each cross-referenced to its resolution.

\begin{enumerate}[leftmargin=1.6em,itemsep=0.5em]
\item \textbf{The bar is overloaded.} $\Prob(A\given B)$ conditions on an
\emph{event} (Definition~\ref{def:elem-cond}); $\E[X\given\Gcal]$ on a
\emph{$\sigma$-algebra} (Definition~\ref{def:cond-exp}); $\E[X\given Y]$ on
a \emph{random variable} (Chapter~\ref{ch:cond-on-rv}); $\E[X\given Y=y]$ is
a \emph{function evaluated} (Definition~\ref{def:eval-ce}); $p(x\given y)$ a
\emph{kernel/density} (Definition~\ref{def:cond-density}). Identify which
before reading on.

\item \textbf{Random variable vs.\ number.} $\E[X\given Y]$ is a random
variable; $\E[X\given Y=y]$ is a (version of a) deterministic function.
Identities about the latter hold only for $P_Y$-a.e.\ $y$
(Warning~\ref{warn:version}).

\item \textbf{``$=$'' usually means ``$\eqas$''.} Conditional objects are
equivalence classes modulo null sets. A clean-looking equation between
conditional expectations is an almost-sure statement unless a specific
version (e.g.\ a regular kernel) has been fixed.

\item \textbf{Same letter $p$, different functions.} $p(x)$, $p(y)$,
$p(x,y)$, $p(x\given y)$ are distinct functions disambiguated only by their
argument names. Restore subscripts $p_X,p_Y,p_{X,Y},p_{X\given Y}$ when in
doubt.

\item \textbf{Semicolon vs.\ bar.} $p(x;\theta)$ / $p_\theta(x)$
(frequentist index, no measure on $\Theta$) vs.\ $p(x\given\theta)$
(Bayesian, $\theta$ random). Only the latter licenses a posterior
$p(\theta\given x)$.

\item \textbf{Reference measure is implicit.} ``Density'' always means
``density with respect to \emph{some} $\lambda$'' --- usually Lebesgue or
counting. Differential entropy and change-of-variables (normalizing flows!)
are sensitive to this choice; mutual information and KL are not
(Warning~\ref{warn:diffent}).

\item \textbf{Conditioning on a null event needs a family.} A single
$p(\cdot\given Y=y)$ is meaningless detached from the disintegrating
variable $Y$; the Borel--Kolmogorov paradox (Section~\ref{sec:bk-paradox})
is the price of forgetting this.

\item \textbf{KL direction matters.} $\KL(p\|q)$ (maximum likelihood,
mode-covering) and $\KL(q\|p)$ (variational inference, mode-seeking) are
different objectives with different failure modes
(Warning~\ref{warn:diffent}'s box, Remark~\ref{prop:elbo}'s discussion).

\item \textbf{$\nabla$ in which variable?} In $\nabla_x\log p_t(x\given y)$
the gradient is in the \emph{state} $x$ (a score), not in parameters
$\theta$; in $\nabla_\theta\log p_\theta(x)$ it is in parameters (for
gradient ascent). The subscript on $\nabla$ is not decoration.

\item \textbf{Expectation subscripts name a measure.} $\E_{x\sim
p(\cdot\given y)}[\,\cdot\,]$ integrates against the kernel
$p(\cdot\given y)$ at fixed $y$ (a conditional expectation); $\E_{x,y}$
integrates against the joint. Reading off the subscript tells you which
measure and hence which Fubini step is legal.
\end{enumerate}

% =====================================================================
\appendix
\part*{Appendices}
\addcontentsline{toc}{part}{Appendices}
% =====================================================================

\chapter{The theorems relied upon}\label{app:theorems}

For completeness we record, without proof, the three measure-theoretic
results that do all the heavy lifting above, plus the topological hypothesis
under which conditional distributions exist. Standard references are
Durrett, \emph{Probability: Theory and Examples}; Kallenberg,
\emph{Foundations of Modern Probability}; and Çınlar, \emph{Probability and
Stochastics}.

\begin{theorem}[Radon--Nikodym]\label{thm:rn}
Let $\mu,\nu$ be $\sigma$-finite measures on $(\Omega,\Fcal)$ with
$\nu\ll\mu$ (i.e.\ $\mu(A)=0\Rightarrow\nu(A)=0$). Then there is an
$\Fcal$-measurable $f\ge0$, unique up to $\mu$-a.e.\ equality, with
$\nu(A)=\int_A f\dd\mu$ for all $A\in\Fcal$. One writes $f=\dd\nu/\dd\mu$.
\end{theorem}

This is the engine of Theorem~\ref{thm:ce-exists}: conditional expectation
\emph{is} a Radon--Nikodym derivative, taken on the sub-$\sigma$-algebra
$\Gcal$. Conditional densities (Definition~\ref{def:cond-density}) and the
KL integrand $\log(\dd P/\dd Q)$ (Definition~\ref{def:kl}) are likewise
Radon--Nikodym derivatives.

\begin{theorem}[Doob--Dynkin lemma]\label{thm:dd-app}
Let $Y:\Omega\to\Ycal$ and $Z:\Omega\to\R$ be measurable with $\Ycal$
standard Borel. Then $Z$ is $\sigma(Y)$-measurable iff $Z=g\circ Y$ for some
Borel $g:\Ycal\to\R$.
\end{theorem}

This is what turns $\E[X\given Y]$ (a $\sigma(Y)$-measurable random
variable) into a function $g$ of $y$, hence into the symbol $\E[X\given
Y=y]$ and the regression function (Chapter~\ref{ch:cond-on-rv}).

\begin{definition}[Standard Borel space]\label{def:std-borel}
A measurable space $(\Xcal,\Bcal)$ is \emph{standard Borel} if it is Borel
isomorphic to a Borel subset of a Polish (separable, completely
metrizable) space with its Borel $\sigma$-algebra. Equivalently, it is
isomorphic to one of: a finite or countable discrete set, or $\R$ (the
Borel isomorphism theorem). All of $\R^d$, separable Banach and Hilbert
spaces, compact metric spaces, $\{0,1\}^\N$, and spaces of probability
measures with the weak topology are standard Borel.
\end{definition}

\begin{theorem}[Existence of regular conditional distributions / disintegration]
\label{thm:disint-app}
Let $X:\Omega\to\Xcal$, $Y:\Omega\to\Ycal$ be random variables with $\Xcal$
standard Borel. Then a regular conditional distribution $\kappa(y,\dd x)$ of
$X$ given $Y$ exists (Definition~\ref{def:rcd}), is $P_Y$-a.e.\ unique, and
disintegrates the joint law as
$P_{X,Y}(\dd x,\dd y)=\kappa(y,\dd x)\,P_Y(\dd y)$
(Theorem~\ref{thm:disintegration}). The standard-Borel hypothesis on
$\Xcal$ is essential; it fails for general measurable spaces.
\end{theorem}

This single hypothesis --- announced as a standing convention in the
Preface and satisfied by every space occurring in practice --- is what makes
the entire conditional vocabulary of machine learning well defined. With it,
$p(x\given y)$, $\pi(a\given s)$, $q_\phi(z\given x)$, and every other bar
in this handbook denotes an object that provably exists.

\chapter{Notation cheat sheet}\label{app:cheatsheet}

\renewcommand{\arraystretch}{1.45}
\begin{longtable}{@{} p{0.30\textwidth} p{0.62\textwidth} @{}}
\toprule
\textbf{Symbol (as printed)} & \textbf{Rigorous meaning / where defined}\\
\midrule
\endfirsthead
\toprule
\textbf{Symbol (as printed)} & \textbf{Rigorous meaning / where defined}\\
\midrule
\endhead
\bottomrule
\endfoot
$\Prob(A\given B)$ &
  Elementary conditional probability, $\Prob(A\cap B)/\Prob(B)$, needs
  $\Prob(B)>0$. Def.~\ref{def:elem-cond}.\\
$\E[X\given\Gcal]$ &
  Conditional expectation: the $\Gcal$-measurable r.v.\ with matching local
  averages; a Radon--Nikodym derivative; the $L^2$-projection onto
  $L^2(\Gcal)$. Def.~\ref{def:cond-exp}.\\
$\E[X\given Y]$ &
  $=\E[X\given\sigma(Y)]$, a $\sigma(Y)$-measurable random variable; equals
  $g(Y)$ by Doob--Dynkin. Ch.~\ref{ch:cond-on-rv}.\\
$\E[X\given Y=y]$ &
  $=g(y)$, a deterministic function defined for $P_Y$-a.e.\ $y$ only.
  Def.~\ref{def:eval-ce}, Warning~\ref{warn:version}.\\
$\Prob(A\given\Gcal)$, $\Prob(A\given Y)$ &
  $\E[\one_A\given\Gcal]$; per-event a.s.-defined random variable.
  Def.~\ref{def:cond-prob-sigma}.\\
$\Var(X\given Y)$ &
  Conditional variance; obeys the law of total variance~\eqref{eq:total-var}.
  Def.~\ref{def:cond-var}.\\
$P_{X\given Y}(\cdot\given y)$, $\kappa(y,\dd x)$ &
  Regular conditional distribution / Markov kernel; exists on standard Borel
  spaces. Defs.~\ref{def:kernel},~\ref{def:rcd}; Thm.~\ref{thm:rcd-exists}.\\
$p(x\given y)$ &
  Conditional density $=p_{X,Y}(x,y)/p_Y(y)$ w.r.t.\ a reference measure;
  a density of the r.c.d. Def.~\ref{def:cond-density}.\\
$p(\theta\given\Dcal)$ &
  Posterior $=$ likelihood $\times$ prior $/$ evidence (Bayes' rule).
  Prop.~\ref{prop:bayes}; \S\ref{sec:posterior}.\\
$X\indep Y\given Z$ &
  Conditional independence; kernel/conditional-expectation/density
  factorization; $\iff \MI(X;Y\given Z)=0$. Def.~\ref{def:ci}.\\
$H(X\given Y)$ &
  Conditional entropy $=\E_Y[H(P_{X\given Y}(\cdot\given Y))]$; expected
  entropy of the r.c.d. Def.~\ref{def:cond-entropy}.\\
$I(X;Y\given Z)$ &
  Conditional mutual information; expected per-fibre KL; $\ge0$, $=0\iff
  X\indep Y\given Z$. Def.~\ref{def:mi}.\\
$\KL(P\|Q)$ &
  Relative entropy $\int\log(\dd P/\dd Q)\dd P$ if $P\ll Q$, else $+\infty$;
  asymmetric. Def.~\ref{def:kl}.\\
$\KL(P_{X\given Y}\|Q_{X\given Y}\,|\,P_Y)$ &
  Conditional / expected KL; the per-step diffusion loss.
  Def.~\ref{def:cond-kl}.\\
$q_\phi(z\given x)$ &
  Variational posterior / encoder kernel; approximates $p_\theta(z\given x)$.
  \S\ref{sec:elbo}.\\
$\E_{q_\phi(z\given x)}[\cdot]$ &
  Integral against the kernel $q_\phi(\cdot\given x)$ at fixed $x$; a
  conditional expectation under $q$. Eq.~\eqref{eq:ce-via-rcd}.\\
$\prod_t p(x_t\given x_{<t})$ &
  Chain rule of probability (an identity, by iterated disintegration), not a
  modeling assumption. \S\ref{sec:autoreg}.\\
$\prod_i p(x_i\given x_{\mathrm{pa}(i)})$ &
  Bayesian-network factorization over a DAG (a Markov assumption).
  Def.~\ref{def:bn}; Thm.~\ref{thm:bn-markov}.\\
$q(x_t\given x_{t-1})$, $p_\theta(x_{t-1}\given x_t)$ &
  Forward / reverse diffusion transition kernels (Markov chain).
  \S\ref{sec:markov-diffusion}.\\
$\nabla_x\log p_t(x)$ &
  Stein score (gradient in state, not parameters); normalizer-free.
  \S\ref{sec:markov-diffusion}.\\
$\nabla_x\log p_t(x\given y)$ &
  Conditional score $=\nabla_x\log p_t(x)+\nabla_x\log p_t(y\given x)$
  (Bayes). Prop.~\ref{prop:guidance}.\\
$\pi(a\given s)$ &
  Policy: a Markov kernel from states to actions. \S\ref{sec:rl}.\\
$V^\pi(s)=\E_\pi[\sum_t\gamma^t R_t\given S_0=s]$ &
  Value: evaluated conditional expectation of discounted return; Bellman $=$
  tower rule. Eqs.~\eqref{eq:value},~\eqref{eq:bellman}.\\
$\E[X_{t+1}\given\Fcal_t]$ &
  Conditioning on a filtration $\sigma$-algebra; martingale if $\eqas X_t$.
  \S\ref{ch:pitfalls} entry, Def.~\ref{def:cond-exp}.\\
\end{longtable}

\vspace{0.5em}
\noindent\emph{One rule to carry away:} a bar is always one of three things
--- a conditional expectation (a projection / Radon--Nikodym derivative on
a $\sigma$-algebra), a Markov kernel (a measurably-indexed family of
measures), or a conditional density (the Radon--Nikodym derivative of that
kernel against a reference measure). Decide which, and the notation stops
being ambiguous.

\end{document}
